\section{Introduction and the main result}\label{sec:introduction}
In the proper hat-guessing game, the players occupy the vertices of a graph.
Adjacent players must wear different colors. Each player sees the colors
of its neighbors and guesses its own color; all guesses are made
simultaneously, and the players win if at least one guess is correct.
A winning strategy must work for every permitted assignment of hats.
Adriaensen et al.\ introduced this variant and the parameter $\HG(G)$
in~\cite{proper}. Table~1 of the inspected version~4 gives
$4\leq\HG(C_5)\leq5$. We settle this particular alternative by excluding
five colors and recalling an explicit four-color strategy.

Related developments include general upper bounds in terms of order and
maximum degree~\cite{kakatelis}, and coverability methods and exact results
for two-spine book graphs~\cite{zhai}. Here we focus on the five-cycle
and the alternative recorded in Table~1 of the cited version of~\cite{proper}.

\begin{theorem}[Main theorem; computer-assisted]\label{thm:main}
The proper hat-guessing number of the five-cycle is
\[
\HG(C_5)=4.
\]
\end{theorem}

The players may use different rules and distinguish left from right.
Their rules need not be invariant under any permutation of vertices or
colors. Nonadjacent hats may agree, and a permitted configuration need
not use all five colors. Thus the task is to exclude the entire local
strategy space, rather than a symmetric family of candidate rules.

\subsection{The idea of the proof}
There are five strategy tables, each with $5^2=25$ ordered observations.
A guess equal to a neighbor's color is always wrong and can be replaced
without losing any previously correct guess. After this normalization,
each of the 125 table entries contributes to exactly one color's count.
Some color $z$ is therefore guessed on at most 25 entries. Write $S(z)$
for this count; it counts table entries, not successful configurations.

To bound $S(z)$ below, consider configurations in which $z$ occurs twice.
The two players wearing $z$ are nonadjacent. Once the answers of the other
three players are fixed, the pairs of colors they miss form a small
bipartite graph. Its vertex covers describe which entries of the two
remaining tables must guess $z$. Every relevant slice needs at least two
such entries, giving a pairwise bound. Classifying when equality is
possible rules out total counts 15 and 16, and yields $S(z)\geq17$.
This is the combinatorial part of the argument.

The remaining work is finite, but the reduction still needs justification.
Vertex relabeling gives 1413 representatives of the necessary five-player
count vectors with $17\leq S(z)\leq25$. These vectors fix only counts,
not the positions of entries or the remaining answers. Each representative
becomes a Boolean formula describing all still-possible winning strategies.
A direct refutation suffices for 1269 representatives. The other 144 are
refined into 51 color-equality patterns each, giving
\[
1269+144\cdot51=8613
\]
terminal formulas. Each has a checked refutation. The choice of which
parents to refine affects computational cost, not the required coverage.

\begin{figure}[htbp]
\centering
\begin{tikzpicture}[
 stage/.style={draw,rounded corners=2pt,align=center,text width=11.8cm,
                inner sep=5pt,font=\small},
 proofarrow/.style={-{Latex[length=2mm]},thick},node distance=7mm]
\node[stage] (strategy) {Assume a five-color winning strategy exists.};
\node[stage,below=of strategy] (count)
 {Normalize, select and relabel a color; choose a count representative\\
 $17\leq S(5)\leq25$ and $a\in\calR_{S(5)}$.};
\node[stage,below=of count] (formula)
 {Choose a direct parent or a color-equivalent diagonal subcase;\\
 extend its strategy assignment to the actual terminal CNF.};
\node[stage,below=of formula] (conflict)
 {That CNF would be satisfiable, but its checked RUP refutation\\
 derives the empty clause. Contradiction.};
\draw[proofarrow] (strategy)--(count);
\draw[proofarrow] (count)--(formula);
\draw[proofarrow] (formula)--(conflict);
\end{tikzpicture}
\caption{The route from an arbitrary strategy to a contradiction.
The reductions preserve a hypothetical winning strategy; the input audit
checks that the generated clauses preserve its satisfying extension.}
\label{fig:proof-route}
\end{figure}

\subsection{What has to be checked}
The proof has three obligations. Keeping them separate makes clear both
what the computer establishes and why that establishes a result about hats.
\begin{enumerate}
\item \emph{Coverage.} Every hypothetical winning strategy has an
 equivalent representative in a listed branch. The mathematical
 reductions are proved in Sections~\ref{sec:preliminaries}--\ref{sec:classification}.
\item \emph{Preservation.} The representative extends to a satisfying
 assignment of the actual branch input. Section~\ref{sec:encoding}
 proves the logical criterion and explains the exhaustive check of the
 generated auxiliary components.
\item \emph{Exclusion.} Every terminal input has a valid refutation.
 Section~\ref{sec:certificates} proves the soundness of the inference
 rule and states the finite assertion verified on the saved files.
\end{enumerate}
A solver's answer alone cannot establish these three facts. In particular,
a refutation of an incorrectly restricted formula would not exclude all
strategies. Section~\ref{sec:conclusion} assembles the three obligations
into the proof of Theorem~\ref{thm:main}.

The switch from a local argument to the full encoding is necessary for
this approach. Appendix~\ref{app:witness} exhibits a strategy with
$S(5)=17$ that wins on all configurations with two hats of color 5, yet
loses on configurations with zero or one such hat. The full formula
therefore includes all 1020 proper five-colorings, not just the
240 configurations used to obtain the combinatorial lower bound.

\subsection{Evidence and reading guide}
The complete computational audit reported here is the run of
27 October 2026: a new input auditor checked every parent and child
formula, and the frozen \texttt{FreshRup} checker replayed every adopted
terminal proof. Appendix~\ref{app:reproduction} gives the versioned
material list and executable reproduction steps;
Appendix~\ref{app:audit} identifies the run records, file fingerprints,
implementation tests and their precise scope. It distinguishes this
full replay from the smaller checks performed while packaging the evidence.

A reader interested primarily in the mathematics can follow
Sections~\ref{sec:preliminaries}--\ref{sec:encoding}, the RUP soundness
lemma, and the final deduction, treating the explicitly stated finite
verification as the computational premise. A reproducer should then
check the actual input files and complete terminal set, using the
appendix rather than saved success summaries as a substitute for running
the checks.

Separate implementations were used within the same AI-assisted workflow.
They are not independent human peer reviewers, and the checker has not
been formally verified. This paper makes no claim of an exhaustive
novelty search: its comparison with~\cite{proper} refers specifically to
the cited version. The evidence is staged for coordinated public release with this
preprint, as described at the end of Section~\ref{sec:conclusion}.
